\documentclass[11pt]{article}
\usepackage{fontspec}
\usepackage[french]{babel}
\usepackage{amsmath}
\usepackage{amssymb}
\usepackage{amsthm}
\usepackage{mathrsfs}
\usepackage{tikz-cd}
\usepackage{enumitem}

\newcommand{\printedpage}[1]{\marginpar{\footnotesize p.~#1}}
\newcommand{\msp}[1]{\marginpar{\footnotesize [#1]}}
\newcommand{\D}{\mathbf{D}}

\begin{document}

\printedpage{39}
% head: A. Grothendieck : Dérivateurs, Ch. II, Édition des Universités de Montpellier II et Paris VII

$I$ satisfaisant la condition a$''$ (autoduale), la condition b$'$ (elle aussi autoduale) est aussi satisfaite (par cas $4^o$ ci-dessus). Je dis donc qu'il existe $I$ et $i$ dans $I$, tels que
\[
\beta_{i,*} : (i\backslash I)^\wedge \longrightarrow I^\wedge
\]
ne soit pas conservatif. (\textbf{NB} Par contre $\alpha_{i,!} : (I/i)^\wedge \longrightarrow I^\wedge$ est toujours conservatif, ce qui redonne le cas $4^o$ comme cas particulier de $3^o$). Comme $\beta_i : i\backslash I \longrightarrow I$ est $W_0$-propre (et même $W_\omega$-propre !), $\beta_{i,*}$ `se calcule fibre par fibre', donc on a
\msp{page 64}
\[
\beta_{i,*}(G)(j) \simeq \prod_{\alpha\in\mathrm{Hom}(i,j)} G(\alpha),
\]
la fibre en question en $j$ étant discrète. Prenons
\[
I = 0 \underset{\beta}{\overset{\alpha}{\rightrightarrows}} 1, \qquad\quad \text{avec } \alpha \neq \beta,
\]
alors $0\backslash I$ est $\Psi = {}$\begin{tikzcd}[ampersand replacement=\&,row sep=0.1em, column sep=1.2em, baseline=(\tikzcdmatrixname-2-1.base)]
 \& 1 \\
0 \ar[ur, "\alpha'"] \ar[dr, "\beta'"'] \& \\
 \& 1'
\end{tikzcd}, et $\Psi^\wedge$ est la catégorie des diagrammes d'ensembles $G = {}$\begin{tikzcd}[ampersand replacement=\&,row sep=0.1em, column sep=1.2em, baseline=(\tikzcdmatrixname-2-1.base)]
 \& G_1 \ar[dl, "\lambda"'] \\
G_0 \& \\
 \& G_{1'} \ar[ul, "\mu"]
\end{tikzcd}. Le foncteur $\beta_*$ déduit de $\beta : i\backslash I \longrightarrow I$ transforme ce diagramme en
\[
G_1 \times G_{1'} \underset{\lambda\,\mathrm{pr}_2}{\overset{\lambda\,\mathrm{pr}_1}{\rightrightarrows}} G_0.
\]
\texttt{[plutôt} $\mu\,\mathrm{pr}_2$\texttt{]}. Si on se borne aux diagrammes $G$ pour lesquels $G_{1'} = \emptyset$, $\beta_*$ sur ceux-ci donne
\[
\emptyset \rightrightarrows G_0,
\]
donc la restriction de $\beta_*$ à cette sous-catégorie (équivalente à $(\boldsymbol{\Delta}^1)^\wedge$) n'est pas fidèle, ni conservative, ne tenant compte que de la composante $G_0$, et non pas de $G_1$.

\msp{page 63bis}

\textbf{Definition 4.6.} \textit{On dit que $u : I' \longrightarrow I$ dans} Diag \textit{est} $\D$-biasphérique \textit{si elle satisfait les conditions équivalentes b, b$'$, b$''$ de 4.4,} $\D$-connexe \textit{si elle satisfait la condition a, ou encore (sous réserve d'existence de $u_!$ resp. de $u_*$) les conditions équivalentes a$'$, a$''$.} % [sic] Definition -> Définition, accent missing in print

Ainsi, $\D$-biasphérique implique $\D$-connexe, et la réciproque est vraie quand $I$ satisfait (par rapport à $\D$) la condition $3^o$ de 4.4 (donc aussi, si elle satisfait l'une des conditions plus faibles $2^o$, $1^o$), ou la condition duale de celle-ci (vu que les notions $\D$-biasphérique, $\D$-connexe sont autoduales). À cet égard, notons le

\textbf{Corollaire 4.7.} \textit{Supposons que $\mathcal{A} = \D(e)$ satisfasse la condition suivante : le foncteur}
\[
\left\{
\begin{array}{lcccl}
\mathcal{A} & \times & \mathcal{A} & \longrightarrow & \mathcal{A} \\[1ex]
(\xi & , & \eta) & \longrightarrow & \xi \amalg \eta
\end{array}
\right.
\]
\textit{est conservatif (resp. fidèle). Alors pour toute flèche étale $u : J \longrightarrow I$ dans} Diag \textit{(et en particulier pour $I/i \longrightarrow I$), le foncteur $u_! : \D(J) \longrightarrow \D(I)$ est conservatif (resp. fidèle)} (\footnote{\textbf{NB} Cette condition est satisfaite si $\forall\, \xi, \eta \in \mathcal{A}$ on a $\mathrm{Hom}(\xi,\eta) \neq \emptyset$, (car on se ramène à ceci : si $u_1$,
% footnote continues at the foot of p. 40
$u_2$ sont deux flèches dans Ens entre ensembles \textit{non vides,} et si $u_1 \times u_2$ est \texttt{[un]} isomorphisme, $u_1$, $u_2$ sont \texttt{[des]} isomorphismes). Cela s'applique en particulier si $\mathcal{A}$ admet un objet neutre, p.ex. si $\mathcal{A}$ est additive.}). \textit{Si cette condition, ou la condition duale,}
\printedpage{40}
% head: A. Grothendieck : Dérivateurs, Ch. II, Édition des Universités de Montpellier II et Paris VII
\msp{page 64bis}
\textit{impliquant les produits $\xi\times\eta$ ($\xi$, $\eta$ dans $\mathcal{A}$), est satisfaite, alors pour toute flèche $u : I' \longrightarrow I$ dans} Diag, \textit{$u$ est $\D$-connexe si et seulement si $u$ est $\D$-biasphérique.}

\textsc{Démonstration.} L'hypothèse de conservativité, ou de fidélité se transporte aux sommes indexées par un ensemble quelconque (sous réserve d'existence), comme on voit en écrivant $\coprod\limits_{\lambda} \xi_\lambda = \xi_{\lambda_0} \amalg \coprod\limits_{\lambda\neq\lambda_0} \xi_\lambda$. L'assertion faite sur $u_!$ résulte alors du calcul de $u_!$ fibre par fibre (qui est valable vu que $u$ est étale donc lisse), en tenant compte que les fibres sont discrètes.

Je ne connais pas d'exemple d'un couple $(\D, u)$ d'un dérivateur $\D$ et d'une flèche $u$ dans Diag, pour lesquels $u$ soit $\D$-connexe sans être $\D$-asphérique. On a donné des conditions sur $\D$, qui assurent équivalence des deux notions. Si $\D$ est de la forme $\D_{\mathcal{M}}$, et si $\mathcal{M}$ satisfait la condition dite dans 4.7 ou sa duale, alors par 4.7 c'est OK. Donc un contrexemple avec un $\D_{\mathcal{M}}$ devrait en tenir compte. Je ne vais pas m'obstiner pour tirer au clair ce qui en est.

\msp{page 65}

\textbf{Corollaire 4.8.} \textit{Une flèche $u$ de} Diag \textit{qui est une $\D$-équivalence universelle est $\D$-biasphérique, et \textup{a fortiori} $\D$-connexe. Inversement, si $u$ est soit $\D$-lisse, soit $\D$-propre, soit une $\D$-fibration faible, et si $u$ est $\D$-biasphérique, c'est une $\D$-équivalence universelle.}

La démonstration est celle du cas particulier 4.3. On a donc un diagramme d'implications :
\[
\begin{tikzcd}[column sep=small, row sep=large]
\text{Cat-équivalence} \ar[dr, Rightarrow] & \D\text{-équivalence universelle } (W_{\mathrm{univ}}) \ar[d, Rightarrow, "{\text{\footnotesize\begin{tabular}{@{}l@{}}implication stricte ; équivalence\\ dans le cas où $u$ est soit $\mathbf{D}$-propre,\\ soit $\mathbf{D}$-lisse, soit une $\mathbf{D}$-fibration faible.\end{tabular}}}"] & \\
 & \D\text{-biasphérique} \ar[dl, Rightarrow] \ar[d, Rightarrow, "{\text{\footnotesize\begin{tabular}{@{}l@{}}équivalence\\ dans tous les\\ cas connus\end{tabular}}}"] \ar[dr, Rightarrow] & \\
\D\text{-asphérique} \ar[dr, Rightarrow] & \D\text{-connexe} \ar[d, Rightarrow] & \D\text{-coasphérique} \ar[dl, Rightarrow] \\
 & \D\text{-équivalences } (W) &
\end{tikzcd}
\]
(\footnote{\textbf{NB} Ces sept classes de morphismes dans Cat sont stables par composition sauf la notion de $\D$-connexité, toutes ces notions ne dépendent que de l'ensemble $W = W_{\D}$ des $\D$-équivalences.}). Rappelons que pour $u : I' \longrightarrow I$ donné
\begin{enumerate}[leftmargin=*, labelindent=1.5em]
\item[a)] $u$ est $\D$-asphérique si et seulement si $\forall\, F$ \textit{constant} dans $\D(I)$
\[
F \xrightarrow[\sim]{\text{adj.}} u_*u^*F \qquad\quad \text{est \texttt{[un]} isomorphisme.} \tag*{$(\ast)$}
\]
\printedpage{41}
% head: A. Grothendieck : Dérivateurs, Ch. II, Édition des Universités de Montpellier II et Paris VII
\item[b)] $u$ est $\D$-coasphérique si et seulement si $\forall\, F$ \textit{constant} dans $\D(I)$
\[
F \xleftarrow[\sim]{\text{adj.}} u_!u^*F \qquad\quad \text{est \texttt{[un]} isomorphisme.} \tag*{$(\ast\ast)$}
\]
\item[c)] $u$ est $\D$-connexe si et seulement si pour \textit{tout} $F$ dans $\D(I)$ (pas forcément constant, ni localement constant), $(\ast)$ est \texttt{[un]} isomorphisme, ou ce qui revient au même, pour tout $F$ $(\ast\ast)$ est \texttt{[un]} isomorphisme.
\end{enumerate}

\msp{page 66}

Rappelons que la classe $W_{\mathrm{asph}}$ \texttt{[des]} flèches $W$-asphériques et celle, $\mathrm{Liss}_W$, des flèches $W$-lisses, se déterminent mutuellement, ainsi
\[
\left\{
\begin{array}{lcl}
(f : X \longrightarrow Y) \in W_{\mathrm{asph}} & \Longleftrightarrow & \forall\, (Y' \xrightarrow{\ u\ } Y) \in \mathrm{Liss}_W, \\
 & & X' = X \times_Y Y' \longrightarrow X' \text{ est } \in W. % [sic] X' = X \times_Y Y' \longrightarrow X' -> X' = X \times_Y Y' \longrightarrow Y', the target must be Y' (compare (4.12))
\\
(f : Y' \longrightarrow Y) \in \mathrm{Liss}_W & \Longleftrightarrow & \text{Le foncteur } X \longmapsto X' = X \times_Y Y' \\
 & & \text{de } \mathrm{Cat}/Y \text{ dans } \mathrm{Cat}/Y' \text{ transforme} \\
 & & \text{flèches dans } W_{\mathrm{asph}} \text{ dans \texttt{[plutôt en]}} \\
 & & \text{flèches dans } W_{\mathrm{asph}}.
\end{array}
\right.
\tag*{(4.10)}
\]

Pour toute partie $A$ de Fl(Cat) stable par isomorphismes, on peut associer une partie $\mathbf{L}(A)$ de Fl(Cat), par la relation
\[
\begin{array}{rcl}
(u : Y' \longrightarrow Y) \in \mathbf{L}(A) & \Longleftrightarrow & u^* : \mathrm{Cat}/Y \longrightarrow \mathrm{Cat}/Y' \\
 & & \text{transforme flèches} \\
 & & \in A \text{ en flèches } \in A.
\end{array}
\tag*{(4.11)}
\]

On voit que $\mathbf{L}(A)$ contient les isomorphismes, est stable par composition, et est stable aussi par changement de base (comme c'est le cas de $\mathrm{Liss}_W$). Inversement, à une partie $L$ de Fl(Cat), stable par isomorphismes, on associe une partie $\mathbf{A}_W(L)$ par
\[
\begin{array}{rcl}
(f : X \longrightarrow Y) \in \mathbf{A}_W(L) & \Longleftrightarrow & \forall\, (Y' \longrightarrow Y) \text{ dans } L, \\
 & & X' = X \times_Y Y' \longrightarrow Y' \\
 & & \text{est dans } W.
\end{array}
\tag*{(4.12)}
\]

Si $L$ est stable par changement de base, alors $\mathbf{A}_W(L)$ est stable par composition. Il contient les isomorphismes (car $W$ les contient), est contenu dans $W$ si $L$ contient les identités. Mais
\msp{page 67}
aucune raison qu'il soit stable par changement de base quelconque (et en fait, $W_{\mathrm{asph}}$ ne l'est pas).

(\footnote{Canulé \texttt{[jusqu'à la fin du} \S4\texttt{]}, cf. \S6.}). On a, pour un $A \subseteq W \subseteq \mathrm{Fl(Cat)}$ donné, stable par isomorphismes,
\[
A \subseteq \mathbf{A}_W(\mathbf{L}(A)) \overset{\text{déf}}{=} A^\natural \subseteq W, \qquad \text{\footnotesize\begin{tabular}{@{}l@{}}partie stable par composition,\\ contient \texttt{[les]} isomorphismes,\end{tabular}}
\tag*{(4.13)}
\]
\printedpage{42}
% head: A. Grothendieck : Dérivateurs, Ch. II, Édition des Universités de Montpellier II et Paris VII
et pour toute partie $L$ de Fl(Cat) qui est stable par changement de base, et par composition
\[
L \subseteq \mathbf{L}(\mathbf{A}_W(L)) \overset{\text{déf}}{=} L^\flat, \qquad \text{\footnotesize\begin{tabular}{@{}l@{}}partie de Fl\,Cat stable\\ par composition et par changement\\ de base, contenant les isomorphismes.\end{tabular}}
\tag*{(4.14)}
\]
On a
\[
\left\{
\begin{array}{lcl}
A^{\natural\natural} & = & A^\natural, \quad \text{car } \mathbf{A}_W(L)^\natural = \mathbf{A}_W(L) \\[1ex]
L^{\flat\flat} & = & L^\flat, \quad \text{car } \mathbf{L}(A)^\flat = \mathbf{L}(A).
\end{array}
\right.
\]
Exemples :
\begin{enumerate}[leftmargin=*, labelindent=1em]
\item[$1^o$)] $L = \text{Ét}$, $\mathbf{A}_W(L) = W_{\mathrm{asph}}$, $\mathbf{L}(\mathbf{A}_W(L)) = L^\flat = \mathrm{Liss}_W$.
\item[$2^o$)] $L = \text{Coét}$, $\mathbf{A}_W(L) = W_{\mathrm{coas}}$, $\mathbf{L}(\mathbf{A}_W(L)) = L^\flat = \mathrm{Prop}_W$.
\item[$3^o$)] $A = W$, $\mathbf{L}(A) = \mathrm{Fib}_W$, $A^\natural = W$.
\item[$4^o$)] $L = \mathrm{Fl\,Cat}$, $\mathbf{A}_W(L) = W_{\mathrm{univ}}$, $L^\flat = L$.
\end{enumerate}
\msp{page 68 vide}

\clearpage
\printedpage{43}
% head: A. Grothendieck : Dérivateurs, Ch. II, Édition des Universités de Montpellier II et Paris VII
\msp{page 69}

\section*{5\quad Dualité formelle}

Soient $\mathcal{P}$, $\mathcal{Q}$ deux ensembles ordonnés. J'essaye de dégager une notion de `dualité' (formelle) reliant $\mathcal{P}$ et $\mathcal{Q}$. Ce sera une relation entre $\mathcal{P}$ et $\mathcal{Q}$
\[
\mathcal{R} \subseteq \mathcal{P} \times \mathcal{Q}
\tag*{(5.1)}
\]
ayant les propriétés suivantes
\begin{enumerate}[leftmargin=*, labelindent=0.5em]
\item[(D)] Pour toute partie $A$ de $\mathcal{P}$, l'ensemble
\end{enumerate}
\[
A^\cup \overset{\text{déf}}{=} \{ q \in \mathcal{Q} \mid (p,q) \in \mathcal{R} \text{ pour tout } p \in A,\ \text{ i.e. } A \times \{q\} \subseteq \mathbf{R} \} % [sic] \mathbf{R} -> \mathcal{R}, the relation is written \mathcal{R} everywhere else
\]
est une partie ouverte de $\mathcal{Q}$ ayant un plus grand élément $\beta(A) : A^\cup = \mathcal{Q}/\beta(A)$. Dualement, pour toute partie $B$ de $\mathcal{Q}$, l'ensemble
\[
B^\cap \overset{\text{déf}}{=} \{ p \in \mathcal{P} \mid (p,q) \in \mathcal{R} \text{ pour tout } q \in B,\ \text{ i.e. } \{p\} \times B \subseteq \mathcal{R} \}
\]
est une partie ouverte de $\mathcal{P}$ ayant un plus grand élément $\alpha(B)$, i.e. $B^\cap = \mathcal{P}/\alpha(B)$.

Le plus souvent, $\mathcal{P}$ et $\mathcal{Q}$ seront stables par sup quelconques (donc aussi par inf quelconques). Dans ce cas, (D) équivaut, du côté des $A^\cup$ (pour $A \in \mathfrak{P}(\mathcal{P})$) à ceci :
\begin{enumerate}[leftmargin=*, labelindent=0em]
\item[(D$_0$)] Pour tout $p \in \mathcal{P}$, $\mathcal{R}(p) = \{ q \in Q \mid (p,q) \in \mathcal{R} \}$ est une partie ouverte \textit{filtrante} de $\mathcal{Q}$, ou ce qui revient au même, une partie non vide ($\Longleftrightarrow$ contenant le plus petit élément $\emptyset_{\mathcal{Q}}$ de $\mathcal{Q}$) stable par sup de deux éléments. De plus, $\mathcal{R}(p)$ est stable par sup filtrants, ou ce qui revient au même (étant lui-même filtrant), il contient un plus grand élément. Il revient % [sic] Q -> \mathcal{Q}, plain italic Q printed in "q \in Q" instead of the script \mathcal{Q}
\msp{page 70}
au même de dire que $\mathcal{R}(p)$ est un \textit{ouvert} de $\mathcal{Q}$ stable par sup quelconques (ce qui implique qu'il contient $\emptyset_{\mathcal{Q}}$). Dualement, pour tout $q \in \mathcal{Q}$, ${}^{s}\mathcal{R}(q) = \{ p \in \mathcal{P} \mid (p,q) \in \mathcal{R} \}$ est une partie \textit{ouverte} de $\mathcal{P}$, stable par sup quelconques, i.e. une partie ouverte ayant un plus grand élément. % doubt resolved against PDF leaf 15: raised s, transpose relation
\end{enumerate}

En d'autres termes, quand on postule la condition (D) pour $A$ resp. $B$ réduit à un élément, alors elle reste vérifiée pour toute partie. Cela provient des formules
\[
\left\{
\begin{array}{lcl}
(\bigcup_i A_i)^\cup & = & \bigcap_i A_i^\cup \\[1ex]
(\bigcup_j B_j)^\cap & = & \bigcap_j B_j^\cap,
\end{array}
\right.
\tag*{(5.2)}
\]
et du fait que l'intersection de parties ouvertes de la forme $\mathcal{P}/\alpha_i$ dans $\mathcal{P}$ est de la forme $\mathcal{P}/\alpha$, où $\alpha = \inf_i \alpha_i$, et de même dans $\mathcal{Q}$. Ainsi on trouve
\[
\left\{
\begin{array}{lcl}
\beta(\bigcup_i A_i) & = & \inf_i \beta(A_i) \\[1ex]
\alpha(\bigcup_j B_j) & = & \inf_j \alpha(B_j).
\end{array}
\right.
\tag*{(5.3)}
\]

\end{document}
